% !TITLE 拟行路难·其四
% !AUTHOR 鲍照
% !DYNASTY 南北朝
% !GENRE 七言古诗
% !ORDER 3
% !EDITION 《先秦汉魏晋南北朝诗》，逯钦立辑校，中华书局
% !EDITION_STATUS pending

\begin{poem}
泻\textsuperscript{1}水置平地，各自东西南北流。
人生亦有命，安能行叹复坐愁？
酌\textsuperscript{2}酒以自宽，举杯断绝\textsuperscript{3}歌路难。
心非木石岂无感？吞声踯躅\textsuperscript{4}不敢言。
\end{poem}

\begin{annotations}
\notetext{1}{泻：泼洒，倾倒。}
\notetext{2}{酌：酌（zhuó），倒酒，斟酒。}
\notetext{3}{断绝：断开、止息。这里指歌声戛然而止，唱不下去。}
\notetext{4}{踯躅：踯躅（zhí zhú），徘徊，走来走去。形容内心情感剧烈冲突、难以决断时的身体动作。}
\end{annotations}

\begin{appreciation}
《拟行路难》共十八首，是鲍照的代表作，这是第四首。它写的是一口咽不下又吐不出的气，中国诗里写“憋屈”最透彻的篇章之一。

“泻水置平地，各自东西南北流”——把水泼在平地上，水便东西南北各流一方。景象平常，道理却重：水往哪里流，取决于地势的高低，与水的意愿无关。人往哪里走，也往往不取决于才具而取决于出身。鲍照身处的南朝，门阀把持要津，世家子弟平流直上寒门之子纵有天大的本事也难进身，他本人正是寒门出身。这个比喻之所以扎人因为他写的是自己的处境。

“人生亦有命，安能行叹复坐愁”——人生自有命数，何必走着也叹、坐着也愁。话说得通达底下却全是一股不服的硬气；真肯认命的人，写不出这样的诗。“酌酒以自宽，举杯断绝歌路难”，斟酒来宽慰自己，酒杯举起来，那支《行路难》却唱不下去了。想唱，却越唱越堵，半句就噎住。

“心非木石岂无感？吞声踯躅不敢言”——心又不是木头石头，怎么会没有感触。可到底还是把声音咽了回去，原地徘徊，不敢开口。“吞声”是把话吞进肚里，“踯躅”是来回打转；两个字眼写尽了一个满是块垒连申诉都不敢的人的窘态。

屈原在《天问》里敢把一百七十多个问题一齐抛给上天，那是他自以为有资格问；鲍照连问的资格都没有，时代不同，人的腰杆也不同。这个题目后来被李白接过去，写下“长风破浪会有时，直挂云帆济沧海”，鲍照咽下去的话，数百年后有人替他吼了出来。杜甫评李白“俊逸鲍参军”，把这份七言里的俊逸认作鲍照的家数。
\end{appreciation}
